% ---------------------------------------------------------------------------
% Chapter: plotkit
% ---------------------------------------------------------------------------
\chapter{plotkit: Publication-Quality Plots}\label{chap:plotkit}
% Long monospaced identifiers leave little room for line breaking; allow some
% extra inter-word stretch in this chapter (the group keeps it local).
\begingroup\setlength{\emergencystretch}{3em}% chapter-local

\pkg{toolkit.plotkit} produces the figures that research papers on
reinforcement learning need most often: learning curves aggregated over seeds
with uncertainty bands, grouped bar charts with error bars, histograms and
annotated heatmaps. It depends only on matplotlib, accepts data from NumPy,
pandas, PyTorch, TensorFlow and JAX, and never changes the global matplotlib
configuration unless explicitly asked to.

\begin{figure}[htbp]
  \centering
  \includegraphics[width=\textwidth]{plotkit_gallery.png}
  \caption{A three-panel summary figure made with \pkg{plotkit}
    (\code{examples/plotkit\_gallery.py}, synthetic data): (a) learning curves
    of three methods over eight seeds with 95\,\% confidence bands and EMA
    smoothing, (b) grouped bars with 95\,\% confidence intervals, (c) an
    annotated heatmap of a hyperparameter sweep. Each method keeps its
    Okabe-Ito colour in every panel.}
  \label{fig:plotkit-gallery}
\end{figure}

\section{Overview}\label{sec:plotkit-overview}

\begin{table}[htbp]
  \centering
  \small
  \caption{Public API of \pkg{toolkit.plotkit}.}
  \label{tab:plotkit-api}
  \begin{tabularx}{\textwidth}{@{}>{\raggedright\arraybackslash}p{4.4cm}>{\raggedright\arraybackslash}X@{}}
    \toprule
    Name & Purpose \\
    \midrule
    \code{plot\_shadow\_curve} & Mean curves with a shaded band (std, SEM, 95\,\% CI, min-max or explicit) and optional smoothing. \\
    \code{plot\_learning\_curves} & \code{\{method: runs\}} learning curves aggregated over seeds; runs may have different lengths. \\
    \code{plot\_line}, \code{plot\_scatter} & Line and scatter plots of one or several series. \\
    \code{plot\_bar} & Bar charts; several series are grouped side by side, with error bars and value labels. \\
    \code{plot\_histogram} & Overlaid histograms on common bins. \\
    \code{plot\_heatmap}, \code{plot\_gray\_scale} & Annotated heatmaps with colorbar (seaborn-compatible arguments, matplotlib only). \\
    \code{save\_figure} & Save a figure to several formats at once. \\
    \code{style\_context}, \code{set\_research\_style}, \code{reset\_style}, \code{STYLE\_PRESETS} & Style presets applied locally or globally. \\
    \code{get\_palette}, \code{RESEARCH\_COLORS}, \code{OKABE\_ITO\_COLORS} (and the \code{*\_COLOR\_LIST} variants) & Colour palettes, including a colour-blind-safe one. \\
    \bottomrule
  \end{tabularx}
\end{table}

All plot functions share the same behaviour:
\begin{itemize}
  \item they draw into the \code{ax} argument when it is given and otherwise
    create a new figure of size \code{figsize};
  \item they apply the \code{style} preset \emph{locally}, inside a
    \code{matplotlib.rc\_context}, so the caller's \code{rcParams} are
    unchanged afterwards;
  \item they return the \code{matplotlib.axes.Axes} they drew on, which can be
    customised further with ordinary matplotlib calls.
\end{itemize}

A non-default style also restyles the axes it draws on (spines, tick labels,
axis-label fonts), so that axes created outside the style context match the
preset. Pass \code{style="default"} to leave an existing axes exactly as it
is.

\begin{pycode}
import numpy as np
from toolkit.plotkit import plot_shadow_curve, save_figure

rng = np.random.default_rng(0)
returns = rng.normal(size=(5, 100)).cumsum(axis=1)     # 5 seeds x 100 steps
ax = plot_shadow_curve(returns, band="ci95", smoothing=10, labels="Agent",
                       xlabel="Episode", ylabel="Return")
save_figure(ax, "renders/first_curve", formats=("pdf", "png"))
\end{pycode}

\section{Input rules}\label{sec:plotkit-inputs}

Every function converts its data inputs with the same rules:

\begin{itemize}
  \item A 1-D array-like of numbers, including a flat Python list of scalars,
    is \emph{one} series.
  \item A list or tuple of array-likes is \emph{several} series, one per
    element; the elements may have different lengths.
  \item A 2-D array is read by the function: \code{plot\_shadow\_curve} reads
    it as \emph{samples of one curve} along \code{axis} (default: rows are
    samples such as seeds, columns are steps); \code{plot\_line},
    \code{plot\_scatter}, \code{plot\_bar} and \code{plot\_histogram} read it
    as \code{(n\_series, n\_points)}.
  \item A pandas \code{DataFrame} contributes one series per column.
  \item pandas \code{Series}, PyTorch tensors (on any device, with or without
    \code{requires\_grad}), TensorFlow eager tensors and JAX arrays are
    converted to NumPy automatically. Masked arrays become NaN.
  \item A single \code{x} is shared by all series; a list of arrays (or a 2-D
    array) gives one \code{x} per series. When \code{x} is omitted, each series
    is plotted against \code{0, \ldots, n - 1}.
\end{itemize}

Length mismatches raise a \code{ValueError} that states both shapes and, when
a 2-D input looks transposed, suggests transposing it. To pass a nested list
as a single 2-D sample matrix to \code{plot\_shadow\_curve}, convert it with
\code{np.asarray} first; otherwise it is read as several curves.

\section{Common parameters}\label{sec:plotkit-common}

The parameters in Table~\ref{tab:plotkit-common} have the same meaning in all
plot functions (the heatmap functions have no \code{labels}, \code{colors},
\code{legend} or \code{grid}).

\begin{table}[htbp]
  \centering
  \small
  \caption{Parameters shared by the plot functions.}
  \label{tab:plotkit-common}
  \begin{tabularx}{\textwidth}{@{}>{\raggedright\arraybackslash}p{2.6cm}>{\raggedright\arraybackslash}p{2.5cm}>{\raggedright\arraybackslash}X@{}}
    \toprule
    Parameter & Default & Meaning \\
    \midrule
    \code{labels} & \code{None} & Legend labels, one per series; default \code{"Curve 1"}, \code{"Series 1"}, \ldots. A wrong count gives a warning. \\
    \code{colors} & \code{None} & \code{None} (the style's colour cycle), one colour, a list of colours (cycled), a palette name (\code{"okabe\_ito"}, \code{"research"}, \ldots) or a matplotlib colormap name. \\
    \code{ax} & \code{None} & Axes to draw on; a new figure is created when \code{None}. \\
    \code{figsize} & \code{(10, 6)} & Size in inches of a new figure (\code{(8, 6)} for heatmaps). \\
    \code{title}, \code{xlabel}, \code{ylabel} & \code{None} & Axes title and axis labels. \\
    \code{legend} & \code{True} & Show a legend when there are several series or labels were given; a string sets the location, e.g. \code{"lower right"}. \\
    \code{grid} & \code{True} & Draw a major grid (only along the value axis for bars and histograms). \\
    \code{style} & \code{"research"} & Style preset applied locally (Section~\ref{sec:plotkit-styles}), a matplotlib style name, a dict of \code{rcParams}, or \code{None}. \\
    \code{**kwargs} & & Passed to the underlying matplotlib call (\code{plot}, \code{scatter}, \code{bar}, \code{hist}, \code{imshow}); \code{color} and \code{label} are accepted as aliases of \code{colors} and \code{labels}. \\
    \bottomrule
  \end{tabularx}
\end{table}

\section{Curves with uncertainty bands}\label{sec:plotkit-curves}

\subsection{plot\_shadow\_curve}

% norun
\begin{pycode}
plot_shadow_curve(y, x=None, y_std=None, labels=None, colors=None, alpha=0.2,
                  ax=None, axis=0, figsize=(10, 6), title=None, xlabel=None,
                  ylabel=None, legend=True, grid=True, style="research",
                  legend_labels=None, x_tick_labels=None, y_tick_labels=None,
                  *, band="std", smoothing=None, **kwargs) -> Axes
\end{pycode}

Each element of \code{y} gives one curve. A 2-D element holds samples of the
curve (for example one row per seed) and is drawn as its mean with a band; a
1-D element is drawn as it is, with a band only when \code{y\_std} is given.

\begin{table}[htbp]
  \centering
  \small
  \caption{Parameters specific to \code{plot\_shadow\_curve}.}
  \label{tab:plotkit-shadow}
  \begin{tabularx}{\textwidth}{@{}>{\raggedright\arraybackslash}p{2.8cm}>{\raggedright\arraybackslash}p{1.6cm}>{\raggedright\arraybackslash}X@{}}
    \toprule
    Parameter & Default & Meaning \\
    \midrule
    \code{y} & required & One or several curves (see the input rules). \\
    \code{x} & \code{None} & Shared x values or one array per curve. \\
    \code{y\_std} & \code{None} & Explicit half-width of the band: a scalar, one array for all curves, or one entry per curve (entries may be \code{None}). Overrides \code{band}. \\
    \code{alpha} & \code{0.2} & Band opacity; \code{0} hides the band. \\
    \code{axis} & \code{0} & Sample axis of 2-D inputs: \code{0} means rows are samples and columns are steps; use \code{1} for the transposed layout. \\
    \code{band} & \code{"std"} & \code{"std"}, \code{"sem"}, \code{"ci95"}, \code{"minmax"} or \code{None} (see below). Keyword only. \\
    \code{smoothing} & \code{None} & \code{int}: trailing moving-average window; \code{float} in $(0,1)$: EMA weight. Keyword only. \\
    \code{legend\_labels} & \code{None} & Alias of \code{labels} (takes precedence). \\
    \code{x\_tick\_labels}, \code{y\_tick\_labels} & \code{None} & Custom tick labels: a mapping \code{\{position: label\}}, or a sequence (one label per x value of the first curve, otherwise spread evenly over the data range). \\
    \bottomrule
  \end{tabularx}
\end{table}

\paragraph{Bands.}
Let $y_{i,t}$ be sample $i$ (for example seed $i$) at step $t$, and let
$\mathcal{I}_t$ be the set of samples that are not NaN at step $t$, with
$n_t = |\mathcal{I}_t|$. The drawn curve is the mean
\begin{equation}
  \bar{y}_t = \frac{1}{n_t} \sum_{i \in \mathcal{I}_t} y_{i,t},
\end{equation}
and the band is $[\bar{y}_t - h_t,\, \bar{y}_t + h_t]$ with
\begin{align}
  \texttt{"std"}:\quad  h_t &= \sigma_t = \sqrt{\frac{1}{n_t}\sum_{i \in \mathcal{I}_t} (y_{i,t} - \bar{y}_t)^2}, \\
  \texttt{"sem"}:\quad  h_t &= \frac{s_t}{\sqrt{n_t}}, \qquad
     s_t = \sqrt{\frac{1}{n_t - 1}\sum_{i \in \mathcal{I}_t} (y_{i,t} - \bar{y}_t)^2}, \\
  \texttt{"ci95"}:\quad h_t &= z_{0.975}\,\frac{s_t}{\sqrt{n_t}}, \qquad z_{0.975} \approx 1.96,
\end{align}
while \code{"minmax"} draws $[\min_{i} y_{i,t},\, \max_{i} y_{i,t}]$ over
$\mathcal{I}_t$. The \code{"std"} band uses the population standard deviation;
the SEM-based bands use the sample standard deviation and are undefined (not
drawn) where fewer than two samples exist. The \code{"ci95"} band is the
normal approximation of a 95\,\% confidence interval of the mean; with few
seeds (fewer than about ten) it is narrower than the exact Student-$t$
interval, a fact worth stating in a figure caption. Because the statistics
ignore NaN, runs of different length can be padded with NaN.

\paragraph{Smoothing.}
Smoothing is applied along the step axis \emph{after} the statistics, with the
same filter for the mean and both band edges. An integer $w \geq 1$ selects a
trailing moving average with the same output length; the window shrinks at
the start of the series:
\begin{equation}
  \tilde{y}_t = \frac{1}{|W_t|} \sum_{s \in W_t} y_s, \qquad
  W_t = \{\, s : \max(0, t - w + 1) \le s \le t,\ y_s \text{ finite} \,\}.
\end{equation}
A float $\beta \in (0, 1)$ selects a debiased exponential moving average with
smoothing weight $\beta$, the convention of TensorBoard's smoothing slider
(larger values are smoother):
\begin{equation}
  \tilde{y}_t = \frac{\sum_{s \le t,\ y_s \text{ finite}} \beta^{\,t-s}\, y_s}
                     {\sum_{s \le t,\ y_s \text{ finite}} \beta^{\,t-s}} .
\end{equation}
The normalisation removes the bias towards zero of the plain recursion
$m_t = \beta m_{t-1} + (1-\beta) y_t$, $m_{-1} = 0$; for finite data the
result equals \code{pandas.Series(y).ewm(alpha=1 - beta).mean()}. Non-finite
values are skipped. Note that smoothing the band edges does not turn them into a
confidence band of the smoothed curve; it is a visual aid.

\subsection{plot\_learning\_curves}

% norun
\begin{pycode}
plot_learning_curves(runs, x=None, *, band="ci95", smoothing=None, colors=None,
                     alpha=0.2, ax=None, figsize=(10, 6), title=None,
                     xlabel="Step", ylabel="Return", legend=True, grid=True,
                     style="research", **kwargs) -> Axes
\end{pycode}

\code{runs} maps a method name to its results: a 2-D array
\code{(n\_seeds, n\_steps)}, a 1-D array (a single seed), or a list of 1-D
arrays of possibly different lengths, which are padded with NaN. \code{x} is
either shared or a mapping \code{\{method: x\}}. The function is a thin layer
over \code{plot\_shadow\_curve} with a 95\,\% confidence band by default, and
always shows the legend when \code{legend} is true.

\section{Lines, scatter plots, bars and histograms}\label{sec:plotkit-basic}

\code{plot\_line(x, y, ...)} and \code{plot\_scatter(x, y, ...)} draw one or
several series with the input rules of Section~\ref{sec:plotkit-inputs}.
\code{plot\_scatter} uses \code{s=50} and \code{alpha=0.7} unless they are
given; passing per-point values \code{c=...} (for a colormap)
disables the series colour.

\subsection{plot\_bar}

% norun
\begin{pycode}
plot_bar(x, height, labels=None, colors=None, ax=None, figsize=(10, 6),
         title=None, xlabel=None, ylabel=None, legend=True, grid=True,
         style="research", *, width=0.8, yerr=None, value_labels=False,
         horizontal=False, capsize=3.0, **kwargs) -> Axes
\end{pycode}

\begin{table}[htbp]
  \centering
  \small
  \caption{Parameters specific to \code{plot\_bar}.}
  \label{tab:plotkit-bar}
  \begin{tabularx}{\textwidth}{@{}>{\raggedright\arraybackslash}p{2.6cm}>{\raggedright\arraybackslash}p{1.6cm}>{\raggedright\arraybackslash}X@{}}
    \toprule
    Parameter & Default & Meaning \\
    \midrule
    \code{x} & required & Category labels (strings or numbers), or \code{None} for the indices. Bars are placed at \code{0, \ldots, n\_categories - 1}. \\
    \code{height} & required & One series (1-D), a list of series, or a 2-D array \code{(n\_series, n\_categories)}. Several series are grouped side by side. \\
    \code{width} & \code{0.8} & Total width of a group in category units; each bar is \code{width / n\_series} wide. \\
    \code{yerr} & \code{None} & Error bars: a scalar, one array per series (or one shared array), or a \code{(2, n\_categories)} array of lower and upper errors for a single series. \\
    \code{value\_labels} & \code{False} & Annotate the bars: \code{True} (format \code{".3g"}), a format string such as \code{".1f"}, \code{"\{:.1\%\}"} or \code{"\%.2f"}, or a callable. \\
    \code{horizontal} & \code{False} & Horizontal bars, first category at the top. \\
    \code{capsize} & \code{3.0} & Error-bar cap size in points. \\
    \code{colors} & \code{None} & As usual; for a single series a list with one colour per category colours each bar. \\
    \bottomrule
  \end{tabularx}
\end{table}

\subsection{plot\_histogram}

\code{plot\_histogram(data, bins="auto", ...)} overlays one histogram per
series; its keyword-only arguments are \code{density=False},
\code{alpha=0.45} and \code{histtype="stepfilled"}. The bin
edges are computed once from the pooled data (\code{numpy.histogram\_bin\_edges})
so that all series share the same bins; NaN and infinite values are dropped.
\code{density=True} normalises each histogram to unit area, and the default
y-label is \code{"Density"} or \code{"Count"} accordingly. For the filled
types the fill uses \code{alpha} while the outline stays opaque.

\section{Heatmaps}\label{sec:plotkit-heatmaps}

% norun
\begin{pycode}
plot_heatmap(data, xlabels=None, ylabels=None, cmap="viridis", annot=False,
             ax=None, figsize=(8, 6), title=None, xlabel=None, ylabel=None,
             cbar=True, cbar_label=None, style="research", *, fmt=None,
             vmin=None, vmax=None, center=None, robust=False, mask=None,
             nan_color=None, aspect="auto", square=False, linewidths=None,
             linecolor="white", annot_kws=None, cbar_kws=None, **kwargs) -> Axes
\end{pycode}

\code{plot\_heatmap} draws a 2-D array with row 0 at the top (a 1-D array is
drawn as one row). It is implemented with matplotlib's \code{imshow} and
accepts the most common seaborn arguments, so existing seaborn calls usually
work unchanged.

\begin{table}[htbp]
  \centering
  \small
  \caption{Parameters specific to \code{plot\_heatmap}.}
  \label{tab:plotkit-heatmap}
  \begin{tabularx}{\textwidth}{@{}>{\raggedright\arraybackslash}p{2.8cm}>{\raggedright\arraybackslash}p{1.9cm}>{\raggedright\arraybackslash}X@{}}
    \toprule
    Parameter & Default & Meaning \\
    \midrule
    \code{xlabels}, \code{ylabels} & \code{None} & One tick label per column / row. \code{None} labels cells with their index (up to 30 cells, otherwise automatic integer ticks); \code{False} hides the ticks. Long x labels are rotated. seaborn's \code{xticklabels} / \code{yticklabels} are aliases. \\
    \code{cmap} & \code{"viridis"} & Colormap name or object; use a diverging map (\code{"RdBu\_r"}) with \code{center}. \\
    \code{annot} & \code{False} & \code{True} writes each value into its cell; an array of the same shape writes its entries instead. Text is black or white, whichever contrasts better with the cell colour. \\
    \code{fmt} & \code{None} & Annotation format (\code{".2g"}, \code{"\{:.1\%\}"}, \code{"\%.2f"} or a callable). Default \code{".0f"} for integer-valued data, else \code{".2g"}. \\
    \code{cbar}, \code{cbar\_label} & \code{True}, \code{None} & Colorbar and its label (also \code{cbar\_kws=\{"label": ...\}}). The colorbar is available as \code{ax.images[-1].colorbar}. \\
    \code{vmin}, \code{vmax} & \code{None} & Colour limits (default: data range). \\
    \code{center} & \code{None} & Value mapped to the middle of the colormap; missing limits are chosen symmetrically around it. \\
    \code{robust} & \code{False} & Use the 2nd and 98th percentiles for missing limits. \\
    \code{mask} & \code{None} & Boolean array; cells where it is \code{True} are hidden. \\
    \code{nan\_color} & \code{None} & Colour of NaN and masked cells (default transparent). \\
    \code{aspect}, \code{square} & \code{"auto"}, \code{False} & Cell aspect ratio; \code{square=True} equals \code{aspect="equal"}. \\
    \code{linewidths}, \code{linecolor} & \code{None}, \code{"white"} & Lines between cells; by default 0.5 for matrices of up to 50 cells per side and none for larger ones. \\
    \code{annot\_kws}, \code{cbar\_kws} & \code{None} & Text properties of the annotations (a \code{color} entry disables the automatic contrast colour) and arguments of \code{Figure.colorbar}. \\
    \bottomrule
  \end{tabularx}
\end{table}

\code{plot\_gray\_scale(data, ...)} is \code{plot\_heatmap} with
\code{cmap="gray"} (black is low, white is high), convenient for images,
occupancy grids and attention maps; it accepts all keyword arguments of
\code{plot\_heatmap}.

\section{Styles}\label{sec:plotkit-styles}

A style preset is a set of \code{rcParams} overrides applied on top of the
active configuration. Table~\ref{tab:plotkit-styles} summarises the presets;
\code{STYLE\_PRESETS[name]} is a read-only view of the exact values.

\begin{table}[htbp]
  \centering
  \small
  \caption{Style presets. Font sizes are in points.}
  \label{tab:plotkit-styles}
  \begin{tabular}{@{}lllll@{}}
    \toprule
     & \code{"research"} & \code{"presentation"} & \code{"minimal"} & \code{"default"} \\
    \midrule
    Font family          & serif & sans-serif & sans-serif & unchanged \\
    Base / label size    & 12 / 12 & 16 / 18 & 10 / 10 & unchanged \\
    Title size, weight   & 14, bold & 20, bold & 11, normal & unchanged \\
    Tick / legend size   & 10 / 10 & 14 / 14 & 9 / 9 & unchanged \\
    Line width           & 2.0 & 3.0 & 1.5 & unchanged \\
    Axes line width      & 1.2 & 1.5 & 0.8 (dark grey) & unchanged \\
    Top and right spines & hidden & hidden & hidden & unchanged \\
    Legend frame         & framed & framed, no edge & none & unchanged \\
    Grid opacity         & 0.3 & 0.3 & 0.2 & unchanged \\
    Colour cycle         & \code{research} & \code{research} & \code{research} & unchanged \\
    \bottomrule
  \end{tabular}
\end{table}

\code{"research"} is the paper-ready default of every function.
\code{"presentation"} enlarges everything for slides and posters,
\code{"minimal"} produces small, quiet figures, and \code{"default"} (or
\code{None}) leaves the active configuration untouched. The \code{style}
argument also accepts any name from \code{matplotlib.style.available} and a
dictionary of \code{rcParams}.

Three functions apply styles outside the plot functions:

\begin{pycode}
import matplotlib.pyplot as plt
from toolkit.plotkit import reset_style, set_research_style, style_context

# 1. Locally, for any matplotlib code (rcParams are restored on exit).
with style_context("presentation"):
    fig, ax = plt.subplots()
    ax.plot([0, 1, 2], [0, 1, 4])

# 2. Globally, as an explicit opt-in; also sets savefig.dpi=300 and
#    savefig.bbox="tight". reset_style() restores the previous rcParams.
set_research_style()
fig, ax = plt.subplots()
ax.plot([0, 1, 2], [0, 1, 4])
reset_style()
\end{pycode}

Library code should prefer the \code{style} argument or
\code{style\_context}; \code{set\_research\_style} is meant for the top level
of a plotting script.

\section{Colour palettes}\label{sec:plotkit-palettes}

\code{get\_palette(name="research", n=None)} returns a list of hex colours:

\begin{itemize}
  \item \code{"research"} (alias \code{"tab10"}): matplotlib's ten default
    colours, also available as the dictionary \code{RESEARCH\_COLORS}
    (\code{"blue"}, \code{"orange"}, \ldots) and the list
    \code{RESEARCH\_COLOR\_LIST}. It is the colour cycle of all presets.
  \item \code{"okabe\_ito"} (alias \code{"colorblind"}): the eight colours of
    Okabe and Ito (2008), designed to remain distinguishable for the common
    forms of colour-vision deficiency; available as \code{OKABE\_ITO\_COLORS}
    and \code{OKABE\_ITO\_COLOR\_LIST}.
  \item Any matplotlib colormap name. Qualitative colormaps (\code{"tab20"},
    \code{"Set2"}, \ldots) return their colours; continuous colormaps
    (\code{"viridis"}, \ldots) are sampled at \code{n} evenly spaced points in
    $[0.1, 0.9]$ (eight by default).
\end{itemize}
Palettes are cycled when \code{n} exceeds their length. A palette name can be
passed directly as \code{colors} to every plot function.

\begin{table}[htbp]
  \centering
  \small
  \caption{The Okabe-Ito palette in plotting order.}
  \label{tab:plotkit-okabe}
  \begin{tabular}{@{}lllll@{}}
    \toprule
    Name & Hex & & Name & Hex \\
    \midrule
    \code{blue}           & \code{\#0072B2} & & \code{reddish\_purple} & \code{\#CC79A7} \\
    \code{orange}         & \code{\#E69F00} & & \code{sky\_blue}       & \code{\#56B4E9} \\
    \code{bluish\_green}  & \code{\#009E73} & & \code{yellow}          & \code{\#F0E442} \\
    \code{vermillion}     & \code{\#D55E00} & & \code{black}           & \code{\#000000} \\
    \bottomrule
  \end{tabular}
\end{table}

\paragraph{Colour-blind guidance.}
About one in twelve men has a red-green colour-vision deficiency. Several
pairs of the default \code{tab10} colours, notably its red and green, are hard
to tell apart with protanopia or deuteranopia, so for papers we recommend
\code{colors="okabe\_ito"}. In addition:
\begin{itemize}
  \item use the first six Okabe-Ito colours for curves; yellow has low
    contrast on white and black carries no hue, so they come last;
  \item keep the colour of a method identical in every panel and figure, by
    passing the same list of colours (Section~\ref{sec:plotkit-recipes});
  \item do not rely on colour alone: vary line styles or markers, or label the
    curves directly, so that the figure also works in greyscale print;
  \item for heatmaps use perceptually uniform colormaps (\code{"viridis"},
    \code{"cividis"}) for magnitudes and a diverging map with \code{center}
    for signed quantities; avoid \code{"jet"} and red-green diverging maps.
\end{itemize}

\section{Saving figures}\label{sec:plotkit-saving}

% norun
\begin{pycode}
save_figure(fig_or_ax, path, formats=None, dpi=300, transparent=False,
            **savefig_kwargs) -> list[Path]
\end{pycode}

\code{save\_figure} accepts a \code{Figure}, a \code{SubFigure}, an
\code{Axes} or an array of axes (the whole parent figure is saved), creates
missing parent directories and returns the written paths. With
\code{formats} given, the extension of \code{path} is replaced by (or, if it
has no recognised extension, extended with) each format: with the path
\code{"results/curve"} and \code{formats=("pdf", "png")}, the files
\code{results/curve.pdf} and \code{results/curve.png} are written. Without
\code{formats} the extension of \code{path} is used, or
\code{rcParams["savefig.format"]} if it has none. \code{dpi} applies to raster
formats and to rasterised parts of vector files. The defaults
\code{bbox\_inches="tight"} and \code{pad\_inches=0.05} can be overridden
through \code{**savefig\_kwargs}.

For LaTeX documents, save vector PDF at the final size of the figure (set
\code{figsize} to the column or text width, for example $3.5 \times 2.5$
inches for a two-column paper, and include it without scaling), so that the
font sizes of the style preset are the sizes that appear in print. Some
publishers require TrueType instead of Type~3 fonts in PDF files; set
\code{pdf.fonttype} while saving:

\begin{pycode}
import matplotlib as mpl

with mpl.rc_context({"pdf.fonttype": 42, "ps.fonttype": 42}):
    save_figure(ax, "renders/curve_truetype.pdf")
\end{pycode}

\section{Recipes}\label{sec:plotkit-recipes}

\subsection{Reinforcement-learning training curves}

Episode returns logged for several seeds form a 2-D array
\code{(n\_seeds, n\_episodes)}. Smoothing with an EMA weight between 0.6 and
0.9 makes the trend visible without hiding the variability between seeds:

\begin{pycode}
import numpy as np
from toolkit.plotkit import plot_shadow_curve

rng = np.random.default_rng(1)
episodes = np.arange(500)
returns = 100 * (1 - np.exp(-episodes / 150)) + rng.normal(0, 15, (10, 500))

ax = plot_shadow_curve(returns, x=episodes, band="ci95", smoothing=0.9,
                       labels="PPO", colors="okabe_ito",
                       xlabel="Episode", ylabel="Episode return",
                       title="Training progress (10 seeds, 95% CI)")
\end{pycode}

\subsection{Comparing methods, including unfinished runs}

\code{plot\_learning\_curves} takes one entry per method. Runs of different
length, for example episodes that ended early or seeds that are still
training, are passed as lists and padded with NaN. The following recipe
combines an environment of \pkg{env\_lib} with \pkg{plotkit}: it records the
order parameter $r$ of the Kuramoto model for three coupling strengths and
eight seeds. Runs with strong coupling synchronise and terminate early, so the
curves end at different steps, and the band of the strongest coupling is
computed from the seeds that are still running.

\begin{pycode}
import numpy as np
import env_lib
from toolkit.plotkit import plot_learning_curves

def order_parameter_trace(coupling, seed, steps=300):
    env = env_lib.KuramotoOscillatorEnv(n_oscillators=10, coupling_mode="constant",
                                        coupling_strength=coupling, max_steps=steps)
    obs, info = env.reset(seed=seed)
    trace = [info["order_parameter"]]
    action = np.zeros(env.action_space.shape, dtype=np.float32)  # no control input
    for _ in range(steps):
        obs, reward, terminated, truncated, info = env.step(action)
        trace.append(info["order_parameter"])
        if terminated or truncated:
            break
    env.close()
    return np.asarray(trace)

runs = {f"K = {k}": [order_parameter_trace(k, seed) for seed in range(8)]
        for k in (0.1, 0.2, 0.5)}
ax = plot_learning_curves(runs, band="ci95", colors="okabe_ito",
                          xlabel="Step", ylabel="Order parameter r",
                          legend="lower right")
\end{pycode}

\subsection{Several panels with consistent colours}

Create the axes yourself, pass them as \code{ax}, and give every panel the
same list of colours. Creating the figure inside \code{style\_context} also
applies the preset to anything drawn directly with matplotlib, such as a
figure title.

\begin{pycode}
import matplotlib.pyplot as plt
import numpy as np
from toolkit.plotkit import (get_palette, plot_bar, plot_learning_curves,
                             save_figure, style_context)

rng = np.random.default_rng(0)
methods = ["Baseline", "Method A", "Method B"]
colors = get_palette("okabe_ito", len(methods))
runs = {m: rng.normal(size=(8, 200)).cumsum(axis=1) + 0.1 * i * np.arange(200)
        for i, m in enumerate(methods)}
finals = np.array([r[:, -1].mean() for r in runs.values()])     # one per method

with style_context("research"):
    fig, axes = plt.subplots(1, 2, figsize=(10, 3.8), layout="constrained")
    plot_learning_curves(runs, colors=colors, smoothing=0.8, ax=axes[0],
                         title="(a) Learning curves")
    plot_bar(methods, finals, colors=colors, ax=axes[1], value_labels=".1f",
             title="(b) Final return")          # one colour per bar
save_figure(fig, "renders/comparison", formats=("pdf", "png"))
\end{pycode}

\subsection{Distinguishing series without colour}

Keyword arguments apply to all series of one call. To give each method its
own line style, call the function once per method on the same axes and pass
the colours explicitly:

\begin{pycode}
fig, ax = plt.subplots(figsize=(6, 4))
for method, color, linestyle in zip(methods, colors, ["-", "--", ":"]):
    plot_shadow_curve(runs[method], labels=method, colors=color, ax=ax,
                      band="sem", linestyle=linestyle)
\end{pycode}

\subsection{Bar chart of final performance}

\begin{pycode}
from toolkit.plotkit import plot_bar

methods = ["DQN", "PPO", "SAC", "Ours"]
accuracy = np.array([71.2, 78.5, 80.1, 84.3])
std_err = np.array([1.8, 1.2, 1.5, 1.1])
ax = plot_bar(methods, accuracy, yerr=std_err, value_labels=".1f",
              colors="okabe_ito", ylabel="Success rate (%)")

# Several environments: height has shape (n_methods, n_environments).
scores = rng.uniform(40, 90, size=(3, 4))
ax = plot_bar(["Env 1", "Env 2", "Env 3", "Env 4"], scores,
              labels=["A", "B", "C"], yerr=rng.uniform(1, 4, size=(3, 4)))
\end{pycode}

\subsection{Confusion and correlation matrices}

\begin{pycode}
from toolkit.plotkit import plot_heatmap

classes = ["left", "stay", "right"]
cm = np.array([[48, 5, 2], [4, 40, 6], [1, 7, 52]])
ax = plot_heatmap(cm, xlabels=classes, ylabels=classes, annot=True,
                  cmap="Blues", square=True, xlabel="Predicted",
                  ylabel="True", cbar_label="Count")

# Correlations: diverging colormap centred at zero, upper triangle hidden.
corr = np.corrcoef(rng.normal(size=(6, 200)))
mask = np.triu(np.ones_like(corr, dtype=bool), k=1)
ax = plot_heatmap(corr, annot=True, fmt=".2f", cmap="RdBu_r", center=0,
                  vmin=-1, vmax=1, mask=mask, square=True, cbar_label="r")
\end{pycode}

\subsection{Distributions of returns}

\begin{pycode}
from toolkit.plotkit import plot_histogram

ax = plot_histogram([rng.normal(0, 1, 1000), rng.normal(1.5, 0.7, 1000)],
                    labels=["Policy A", "Policy B"], density=True,
                    colors="okabe_ito", xlabel="Return")
\end{pycode}

\subsection{Plotting tensors and data frames}

Tensors can be passed without conversion; the functions detach them and copy
them to the CPU. A pandas \code{DataFrame} gives one series per column, which
suits logs read with \code{pandas.read\_csv}.

\begin{pycode}
import torch
from toolkit.plotkit import plot_line

losses = torch.rand(3, 100, requires_grad=True).cumsum(dim=1)   # (n_series, n_points)
ax = plot_line(None, losses, labels=["actor", "critic", "entropy"])
\end{pycode}

\section{Command-line gallery}\label{sec:plotkit-cli}

The module is executable and shows every plot type with synthetic data. The
same program is installed as the console script \code{plotkit-gallery}.

\begin{shellcode}
python -m toolkit.plotkit --demo all                    # all eight demos in one window
plotkit-gallery --demo heatmap --style presentation \
    --save renders/plotkit --format png pdf             # renders/plotkit/plotkit_heatmap.*
\end{shellcode}

\begin{table}[htbp]
  \centering
  \small
  \caption{Options of \code{python -m toolkit.plotkit} and \code{plotkit-gallery}.}
  \label{tab:plotkit-cli}
  \begin{tabularx}{\textwidth}{@{}>{\raggedright\arraybackslash}p{3.0cm}>{\raggedright\arraybackslash}p{2.1cm}>{\raggedright\arraybackslash}X@{}}
    \toprule
    Option & Default & Meaning \\
    \midrule
    \code{--demo} & \code{all} & \code{shadow}, \code{learning}, \code{heatmap}, \code{grayscale}, \code{line}, \code{bar}, \code{scatter}, \code{histogram}, or \code{all} for a $2 \times 4$ grid of all demos. \\
    \code{--style} & \code{research} & Style preset. \\
    \code{--save DIR} & none & Save to \code{DIR/plotkit\_<demo>.<format>} instead of showing a window. \\
    \code{--format} & \code{png} & One or more output formats. \\
    \code{--dpi} & \code{150} & Resolution of raster output. \\
    \code{--seed} & \code{0} & Seed of the synthetic data. \\
    \bottomrule
  \end{tabularx}
\end{table}

Without \code{--save} the figure is shown with \code{plt.show()}; when the
active matplotlib backend is non-interactive (for example with
\code{MPLBACKEND=Agg}), the figure is saved to \code{renders/plotkit/}
instead. The script \code{examples/plotkit\_gallery.py} builds the
three-panel figure of Figure~\ref{fig:plotkit-gallery} and is a good starting
point for a paper figure; run it with \code{--help} for its options.

\par\endgroup
